% Supplementary Information of the ERL Letter, section \SIrevision: the analyses added after the
% internal review (PREREG Amendment 18) and the classification of every
% quantitative claim of the Abstract and Results. Every number is a
% \metric{} macro.

All analyses in this section were registered as Amendment 18 of the pre-registration before any of them was computed, and are exploratory unless a table labels them otherwise. No registered verdict is replaced. The registration, its eighteen amendments and its erratum are reproduced verbatim, each with the commit and date that first recorded it, in the Supplementary file \emph{Pre-registration record} (supplementary\_prereg.md).

\subsection*{Classification of the quantitative claims}
Table~\ref{tab:si_claims} labels every quantitative claim of the Abstract and the Results of the Letter. \emph{Confirmatory}: a registered decision rule, or a registered metric reported as registered. \emph{Registered sensitivity}: a variant or check named in the registration before the results. \emph{Exploratory}: not registered before the data or results were seen, or registered in Amendment 18 after the internal review.

\begin{table}[!htbp]
\centering
\scriptsize
\setlength{\tabcolsep}{3pt}
\caption{\textbf{Status of every quantitative claim of the Abstract and Results.} Where, the place in the Letter (A, Abstract; R, Results subsection); the registration items are in the Supplementary file \emph{Pre-registration record}.}
\label{tab:si_claims}
\begin{tabular}{>{\raggedright\arraybackslash}p{7.6cm}>{\raggedright\arraybackslash}p{1.6cm}>{\raggedright\arraybackslash}p{4.6cm}}
\hline
Claim & Where & Status \\
\hline
Fall of the basin level between its smoothed maximum and minimum & A & Descriptive input (no test) \\
Four grids mispredict all four validated sectors; displacement, floors and $D/Q$ there & A, R first & Exploratory (Amendment 18); the counts use the registered rule \\
Four grids mispredict six or seven of seven sectors after the shift; all five, every sector without it & R first & Confirmatory (H1 rule, registered section 9) \\
Bootstrap probabilities of misprediction & R first & Exploratory (Amendment 18) \\
Open shelf without the record-low rule & A, R first & Registered sensitivity (red-team run); its reading is exploratory \\
Displacements in the deltas and Komsomolets Bay & R first & Confirmatory metric; the label ``not validated'' is exploratory \\
Verdicts with per-sector shifts & R first & Registered sensitivity (red-team run) \\
Verdicts on test blocks only and with a halves-only floor & R first & Exploratory (Amendment 18) \\
Whole-metre control, static filter, resolution-matched floor, displacement-fitted shift, error budget & A, R first & Exploratory (Amendment 18; the control's stop rule registered there) \\
Seabed change between epochs (comparisons, binomial test, hysteresis) & R first & Registered check (section 9); the test and hysteresis are exploratory \\
Median area bias after the shift and on the datum chain & A, R second & Confirmatory metric (no decision rule) \\
Area bias against level, out of sample, intersection over union in the band & A, R second & Exploratory (Amendment 18) \\
Deltaic block permutation test & R second & Confirmatory (H1 rule); sector-level test and Holm adjustment exploratory \\
ETOPO 2022 mispredicts one sector, none validated & A, R third & Confirmatory count; validated-sector reading exploratory \\
ETOPO 2022 lake processing, shoreline sources, $D/Q$ by band, GEDI contrast & A, R third & Exploratory (Amendment 18; band edges set after the results) \\
Great Salt Lake: every grid mispredicts both arms; the control passes & A, R fourth & Confirmatory (Amendment 16; registered substitute control) \\
Great Salt Lake: grid sources, verdicts at the control's levels, pooled floor, lidar shares & A, R fourth & Exploratory (Amendment 18) \\
Near-term emergence from the observed 2026 shoreline (composite baseline) and on the shifted grids & A, R fifth & Exploratory (Amendment 18; the composite baseline was added after the registered baseline's flaw was seen) \\
Near-term emergence on the registered common baseline; the running-median 2026 level & S9 only & Registered sensitivity (section 10) \\
Factor between the grids as distributed, reference-level change, already-dry share & R fifth & Exploratory (Amendment 18) \\
Hindcast envelopes, vertical-error simulation and agreement reference, dry against isolated water, pathway by grid variance shares, trend windows & A, R fifth & Exploratory (Amendment 18) \\
Protected-area share with and without one polygon & R fifth & Exploratory (Amendment 18; a diagnostic) \\
Dates at which the primary pathways reach each decline & R fifth & Confirmatory metric (section 10) \\
Registered correction, its transfer test and its verdict; ATL24 check & R fifth & Confirmatory (H3; Amendment 17) \\
H2 verdict, NMAD and RMSE of the elevation model (Discussion) & --- & Confirmatory (H2 rule) \\
H2 split draws, influence of blocks and sectors, coverage intervals (Discussion) & --- & Exploratory (Amendment 18) \\
\hline
\end{tabular}
\end{table}

\subsection*{H1 in the validated sectors}
Table~\ref{tab:si_rev_sectors} gives, for the four grids that mispredict, the per-sector results that the Letter's headline draws on. The moving-block bootstrap resamples levels in blocks of \metric{rev.w.s.boot.block_m}~m (\metric{rev.w.s.boot.n} resamples, seed \metric{rev.w.s.boot.seed}). In the validated sectors the probability that a grid mispredicts all four is \metric{rev.s.item2.count.shape.validated4.majority.p_all_min}--\metric{rev.s.item2.count.shape.validated4.majority.p_all_max}; over all seven sectors it is only \metric{rev.s.item2.count.shape.all7.majority.p_all_min}--\metric{rev.s.item2.count.shape.all7.majority.p_all_max}, because the Komsomolets Bay verdicts are near even. In Komsomolets Bay, \metric{rev.s.item2.dq.shape.gebco2026.komsomolets.n_levels_q_over_10km} of the evaluable levels have floors above ten kilometres, which is why a sector median of \metric{rev.s.item2.dq.shape.gebco2026.komsomolets.d_km}~km (GEBCO\_2026) passes the rule by a single level. Excluding Komsomolets Bay, the same grids mispredict \metric{rev.s.item1.shape.six_no_komsomolets.majority_products_min} of the six remaining sectors. ETOPO 2022 mispredicts none of them, with a bootstrap probability of \metric{rev.w.h1.etopo2022.validated4.p_at_least_one} of mispredicting at least one.

\begin{table}[!htbp]
\centering
\footnotesize
\caption{\textbf{H1 per sector, shape variant, for the four grids that mispredict} (GEBCO\_2026, GEBCO\_2022, ETOPO1 and SRTM15+): the range over the four grids of the sector median displacement $D$, the sector floor $Q$ (median over evaluable levels of the per-level \metric{params.err_pct_hi}th percentile), the range of the median per-level ratio $D/Q$, and the range of the bootstrap probability of misprediction.}
\label{tab:si_rev_sectors}
\begin{tabular}{llllll}
\hline
Sector & Observations & $D$ (km) & $Q$ (km) & $D/Q$ & P(mispredict) \\
\hline
West & validated & \metric{rev.s.headline.west.d_km_min}--\metric{rev.s.headline.west.d_km_max} & \metric{rev.s.headline.west.q_km} & \metric{rev.s.headline.west.dq_median_min}--\metric{rev.s.headline.west.dq_median_max} & \metric{rev.s.headline.west.p_mispredict_min}--\metric{rev.s.headline.west.p_mispredict_max} \\
Interfluve & validated & \metric{rev.s.headline.interfluve.d_km_min}--\metric{rev.s.headline.interfluve.d_km_max} & \metric{rev.s.headline.interfluve.q_km} & \metric{rev.s.headline.interfluve.dq_median_min}--\metric{rev.s.headline.interfluve.dq_median_max} & \metric{rev.s.headline.interfluve.p_mispredict_min}--\metric{rev.s.headline.interfluve.p_mispredict_max} \\
Northeastern coast & validated & \metric{rev.s.headline.northeast.d_km_min}--\metric{rev.s.headline.northeast.d_km_max} & \metric{rev.s.headline.northeast.q_km} & \metric{rev.s.headline.northeast.dq_median_min}--\metric{rev.s.headline.northeast.dq_median_max} & \metric{rev.s.headline.northeast.p_mispredict_min}--\metric{rev.s.headline.northeast.p_mispredict_max} \\
Open shelf & validated & \metric{rev.s.headline.open_shelf.d_km_min}--\metric{rev.s.headline.open_shelf.d_km_max} & \metric{rev.s.headline.open_shelf.q_km} & \metric{rev.s.headline.open_shelf.dq_median_min}--\metric{rev.s.headline.open_shelf.dq_median_max} & \metric{rev.s.headline.open_shelf.p_mispredict_min}--\metric{rev.s.headline.open_shelf.p_mispredict_max} \\
Volga delta & not validated & \metric{rev.s.headline.volga.d_km_min}--\metric{rev.s.headline.volga.d_km_max} & \metric{rev.s.headline.volga.q_km} & \metric{rev.s.headline.volga.dq_median_min}--\metric{rev.s.headline.volga.dq_median_max} & \metric{rev.s.headline.volga.p_mispredict_min}--\metric{rev.s.headline.volga.p_mispredict_max} \\
Ural delta & not validated & \metric{rev.s.headline.ural.d_km_min}--\metric{rev.s.headline.ural.d_km_max} & \metric{rev.s.headline.ural.q_km} & \metric{rev.s.headline.ural.dq_median_min}--\metric{rev.s.headline.ural.dq_median_max} & \metric{rev.s.headline.ural.p_mispredict_min}--\metric{rev.s.headline.ural.p_mispredict_max} \\
Komsomolets Bay & not validated & \metric{rev.w.h1.komsomolets.d_km_min}--\metric{rev.w.h1.komsomolets.d_km_max} & \metric{rev.w.h1.komsomolets.q_km_min} & \metric{rev.w.h1.komsomolets.dq_min}--\metric{rev.w.h1.komsomolets.dq_max} & \metric{rev.w.h1.komsomolets.p_min}--\metric{rev.w.h1.komsomolets.p_max} \\
\hline
\end{tabular}
\end{table}

\emph{The record-low rule.} It is built from E3 bins \metric{rev.s.r2minor8.bins_text} (\metric{rev.s.r2minor8.level_low_m} to \metric{rev.s.r2minor8.level_high_m}~m), \metric{rev.s.r2minor8.n_acquisitions} acquisitions from \metric{rev.s.r2minor8.first} to \metric{rev.s.r2minor8.last} in months \metric{rev.s.r2minor8.months_text}; the lowest audited level is \metric{rev.s.r2minor8.lowest_evaluable_level_m}~m in the interfluve and Ural delta and \metric{rev.s.r2minor8.lowest_evaluable_level_max_m}~m elsewhere. Because the rule is applied to both halves, it forces agreement in the cells it fixes and lowers the floor where it acts. On the open shelf it lowers the floor from \metric{rev.s.r3minor1.open_shelf.floor_noenv_km} to \metric{rev.s.r3minor1.open_shelf.floor_primary_km}~km and decides the verdict (\metric{rev.s.r3minor1.open_shelf.main.shape.n_products} grids mispredict with it, \metric{rev.s.r3minor1.open_shelf.main_noenv.shape.n_products} without); elsewhere it changes the floors by a factor of \metric{rev.s.r3minor1.volga.ratio}--\metric{rev.s.r3minor1.northeast.ratio}. Without the rule the four grids mispredict \metric{rev.w.h1.noenv.validated4.n_min} of the four validated sectors. The envelope's bins include the two lowest audited levels, whose extents are therefore partly set by the rule itself.

\emph{The floor, test blocks and the open-shelf medians.} The registered floor pools, at levels crossed on both limbs, the distances between the E1 shoreline and the later ones; since the E1 halves never produce a distance, the E1 part of the floor is entirely between epochs. With a floor from the halves alone, the GEBCO grids and ETOPO1 are unchanged (\metric{rev.w.h1.halves.gebco_etopo1.n_min} of the seven sectors), SRTM15+ goes to \metric{rev.s.r3m4.halves_floor.shape.srtm15p.n_sectors} and ETOPO 2022 to \metric{rev.s.r3m4.halves_floor.shape.etopo2022.n_sectors}, so the registered floor slightly favoured the grids. On test blocks alone the four grids mispredict \metric{rev.w.h1.test.validated4.n_min} of the four validated sectors; on the open shelf GEBCO\_2026's $D$ is then \metric{rev.s.r3m4.test.shape.gebco2026.open_shelf.d_km}~km, against \metric{rev.s.item2.dq.shape.gebco2026.open_shelf.d_km}~km over all blocks, from few test distances. GEBCO\_2026's open-shelf median is \metric{rev.s.r1minor11.shape.gebco2026.open_shelf.median_observed_km}~km over all \metric{rev.s.r1minor11.shape.gebco2026.open_shelf.n_observed} observed levels and \metric{rev.s.r1minor11.shape.gebco2026.open_shelf.median_evaluable_km}~km over the \metric{rev.s.r1minor11.shape.gebco2026.open_shelf.n_evaluable} evaluable ones; every table now uses the evaluable levels.

\emph{Resolution, quantisation and the systematic error budget.} A rerun of the audit reproduces the registered displacements and floors exactly (largest difference \metric{rev.r.replica.levels_D_max_abs_diff_m}~m). \emph{Control.} The waterline elevation model, rounded to whole metres and aggregated to each grid's native cells, was flooded and audited like a grid (its area-fitted constants \metric{rev.r.control.const.u15_m} to \metric{rev.r.control.const.r15_gebco_m}~m). Its median displacement is \metric{rev.w.r.control.westinter_km_min}--\metric{rev.w.r.control.westinter_km_max}~km on the west and interfluve, \metric{rev.w.r.control.coasts3_km_min}--\metric{rev.w.r.control.coasts3_km_max}~km over the three validated coasts and \metric{rev.w.r.control.openshelf_km_min}--\metric{rev.w.r.control.openshelf_km_max}~km on the open shelf. The registered stop rule does not fire: the grids exceed the control's displacement at most levels in \metric{rev.r.control.stop.n_excess_holds} of \metric{rev.r.control.stop.n_cases} grid-by-validated-sector cases (\metric{rev.r.control.stop.unrounded_alllevels.n_excess_holds} of \metric{rev.r.control.stop.n_cases} against the unrounded control); the \metric{rev.r.control.stop.n_excess_fails} failures are GEBCO\_2022 and SRTM15+ on the open shelf. SRTM15+ also stores most of its band in whole metres. \emph{Static-contour filter.} Removing static channel-bank waterlines changes \metric{rev.r.static.n_changed} verdicts; it changes the floors of the Volga delta (\metric{rev.r.static.floor_primary.volga_km} to \metric{rev.r.static.floor_static.volga_km}~km), the west (\metric{rev.r.static.floor_primary.west_km} to \metric{rev.r.static.floor_static.west_km}~km) and the Ural delta (\metric{rev.r.static.floor_primary.ural_km} to \metric{rev.r.static.floor_static.ural_km}~km). Static shares of the audit cells are \metric{rev.r.static.share_audit_cells.volga_pct}\% (Volga), \metric{rev.r.static.share_audit_cells.west_pct}\% (west), \metric{rev.r.static.share_audit_cells.ural_pct}\% (Ural), \metric{rev.r.static.share_audit_cells.openshelf_pct}\% (open shelf), \metric{rev.r.static.share_audit_cells.interfluve_pct}\% (interfluve), \metric{rev.r.static.share_audit_cells.northeast_pct}\% (northeastern coast) and \metric{rev.r.static.share_audit_cells.komsomolets_pct}\% (Komsomolets Bay). \emph{Resolution-matched floor.} With the observed extents aggregated to each grid's cells, \metric{rev.r.res.n_changed} verdicts change: the GEBCO grids no longer mispredict the open shelf and ETOPO1 no longer the northeastern coast, which leaves \metric{rev.w.r.res.validated4_min}--\metric{rev.w.r.res.validated4_max} validated sectors per grid. \emph{Constants.} The displacement-fitted constant lies \metric{rev.w.r.cD_below_min_m}--\metric{rev.w.r.cD_below_max_m}~m below the area-fitted one and changes \metric{rev.r.constants.n_changed_at_c_D} verdicts; with the best constant chosen separately in each sector, the shorelines still lie beyond the floor in \metric{rev.w.r.bestcase.validated4_min}--\metric{rev.w.r.bestcase.validated4_max} validated sectors per grid (the open shelf, and for GEBCO\_2022 the northeastern coast, can be fitted). \emph{Gradients and error budget.} Median shelf gradients are \metric{rev.w.r.gmedian_min}--\metric{rev.w.r.gmedian_max}~m per km, so the shoreline moves about \metric{rev.w.r.km_per_m_min}--\metric{rev.w.r.km_per_m_max}~km per metre of level. A systematic vertical budget (level uncertainty, the Hydroweb shift error, setup residuals and local departures; standard deviation \metric{rev.w.r.budget.sigma_median_min}--\metric{rev.w.r.budget.sigma_median_max}~m) corresponds to \metric{rev.w.r.budget.horizontal_median_min}--\metric{rev.w.r.budget.horizontal_median_max}~km; added to the floor it changes \metric{rev.r.budget.n_changed} verdicts: all four grids pass the open shelf and GEBCO\_2022 and SRTM15+ the northeastern coast, leaving \metric{rev.w.r.budget.validated4_min}--\metric{rev.w.r.budget.validated4_max} validated sectors per grid; the west and the interfluve still fail for all four. The rough gradient of the elevation model probably overstates slopes, which makes this floor smaller. \emph{Diagnostics.} The share of flood-season scenes rises with level, but does not track the displacement residuals in the deltas (Spearman $\rho$ \metric{rev.r.season.rho_min} to \metric{rev.r.season.rho_max}; no test significant). Converted with the sector gradients, the grids' ICESat-2 dispersion ranks them only weakly like their displacement ($\rho$ \metric{rev.r.nmad.rho_overall.atl08}--\metric{rev.r.nmad.rho_overall.atl06}). Not run, for time: the threshold and vegetation variants through the hindcast, and these checks for ETOPO 2022.

\emph{Water surface and datums.} \emph{Departures.} From the Surface Water and Ocean Topography (SWOT) block means before their removal (\metric{rev.v.dep.all.n_passes} warm-season passes, \metric{rev.v.dep.first_pass} to \metric{rev.v.dep.last_pass}), the water surface lies \metric{rev.w.dep.common_m}~m from the DAHITI level over all blocks; this common value includes reference offsets and is not used on its own. Relative to it the sectors depart by \metric{rev.v.dep.rel_min_m} to \metric{rev.v.dep.rel_max_m}~m, the largest in the Volga delta; Komsomolets Bay has no open-sea block. Shifting each sector's levels by its departure changes \metric{rev.w.v.depaudit.shape_n_changed} of the \metric{rev.w.v.depaudit.n_cases} shape-variant verdicts. \emph{Planar field.} A plane fitted per grid by the registered area criterion lowers the area objective by \metric{rev.w.v.planar.reduction_pct_min}--\metric{rev.w.v.planar.reduction_pct_max}\%, but the median of the sector medians of displacement grows from \metric{rev.w.v.planar.gebco2026.before_km} to \metric{rev.w.v.planar.gebco2026.after_km}~km for GEBCO\_2026 and from \metric{rev.w.v.planar.etopo1.before_km} to \metric{rev.w.v.planar.etopo1.after_km}~km for ETOPO1 (the two grids audited), and every sector stays mispredicted: a smooth datum field fitted on area improves the totals and worsens the shorelines. \emph{Posts.} None of the \metric{datum.bhs.n_posts} posts behind $\delta$ lies on the North Caspian shelf; their offsets depart from $\delta$ by \metric{rev.w.v.posts.delta_dev_m_min} to \metric{rev.w.v.posts.delta_dev_m_max}~m. The four North Caspian posts, not used for $\delta$, depart by \metric{rev.v.posts.lag.minus_delta_m}~m (Lagan), \metric{rev.v.posts.tul.minus_delta_m}~m (Tyuleniy), \metric{rev.v.posts.pesh.minus_delta_m}~m (Peshnoy) and \metric{rev.v.posts.kul.minus_delta_m}~m (Kulaly). \emph{The ICESat-2 constant.} $c_\mathrm{ref}$ is \metric{rev.w.v.cref.c_m_min}--\metric{rev.w.v.cref.c_m_max}~m. The documented terms, the zero-degree term\citep{iers_tn36} and the permanent tide, leave \metric{rev.w.v.cref.resid_documented_m_min}--\metric{rev.w.v.cref.resid_documented_m_max}~m unexplained. Heights above the T/P ellipsoid are about \metric{rev.v.cref.tp_term_min_m}~m larger than above WGS84\citep{is2_compguide_r007}; an undocumented reduction of DAHITI to that ellipsoid would leave \metric{rev.w.v.cref.resid_tp_m_min}--\metric{rev.w.v.cref.resid_tp_m_max}~m. SWOT, which uses no elevation model, also puts the DAHITI level above the water surface, so the offset lies in DAHITI's reference, and the total variant's datum error includes it. \emph{The EGM2008 constant.} NGA's realisation of EGM2008 adds \metric{datum.nga_offset_m}~m to the geoid height\citep{geographiclib_gravity}; ICGEM's evaluation of the same coefficients does not\citep{icgem_faq}. \emph{Constants in the Baltic Height System.} Expressed in BHS-77, the area-fitted constants mean adding \metric{rev.w.v.bhs.k_m_min}--\metric{rev.w.v.bhs.k_m_max}~m to the heights as distributed, and the per-sector constants spread by \metric{rev.w.v.bhs.sector_spread_min_m}--\metric{rev.w.v.bhs.sector_spread_max_m}~m. The charts behind the grids add their own references: ETOPO1's Caspian bathymetry is gridded from the contours of a 1987 navigation map referenced to the local sea level and shifted uniformly\citep{amante2009}, and GEBCO's Caspian grid, from 2010 at least to GEBCO\_2014, is a compilation of Russian hydrographic charts, shifted likewise, with no chart datum named\citep{gebco2014_doc}; the current grid still codes the area as a pre-generated grid\citep{gebco2026}. No document gives the charts' zero in BHS-77, and a compilation of sheets surveyed at different sea levels could leave an offset that varies in space. \emph{Levels.} The fall of \metric{levels.dahiti.smoothed.range_m}~m is the range of the three-record running median; the single records span \metric{rev.v.levels.fall_raw_m}~m, and the annual means fall from \metric{rev.w.levels.annual_1995_bhs_m}~m (1995) to \metric{rev.w.levels.annual_2025_bhs_m}~m (2025) in BHS-77, by \metric{rev.w.levels.annual_fall_m}~m. The 2026 reference level is the January--September mean, \metric{rev.w.levels.janSep_above_annual_m}~m above an annual mean, so its difference from the end-2024 level is a seasonal window, not a rise. DAHITI holds \metric{rev.v.levels.dahiti_records_2024}--\metric{rev.v.levels.dahiti_records_2025} records a year in 2024--2025, and \metric{rev.w.levels.e3_hydroweb_pct}\% of the late-fall scenes below \metric{h1.disp_band.edge_1_m}~m take shifted Hydroweb levels; these enter the error budget. The mean depth of the domain's sea at the 2026 level is \metric{rev.w.v.depth2026_m_min}--\metric{rev.w.v.depth2026_m_max}~m on the shifted grids.

\emph{Area-based claims.} After the shift, the per-level area bias has interquartile ranges of \metric{rev.w.s.iqr.domain_min}--\metric{rev.w.s.iqr.domain_max} percentage points over the domain and \metric{rev.w.s.iqr.validated_min}--\metric{rev.w.s.iqr.validated_max} in the validated sectors. Its Theil--Sen slope against level is \metric{rev.w.s.slope.domain_min}--\metric{rev.w.s.slope.domain_max}\% per metre for the four grids (lower \metric{params.ci_pct}\% bounds \metric{rev.w.s.slope_lo.domain_min}--\metric{rev.w.s.slope_lo.domain_max}) and \metric{rev.w.s.slope.validated_min}--\metric{rev.w.s.slope.validated_max}\% per metre in the validated sectors (lower bounds \metric{rev.w.s.slope_lo.validated_min}--\metric{rev.w.s.slope_lo.validated_max}); for ETOPO 2022 it is \metric{rev.s.item4.area_bias.shape.etopo2022.domain.theil_sen_pct_per_m}\% per metre (\metric{rev.s.item4.area_bias.shape.etopo2022.domain.theil_sen_lo} to \metric{rev.s.item4.area_bias.shape.etopo2022.domain.theil_sen_hi}). Fitted on the training blocks of E1 and E2 only, the shift is \metric{rev.w.s.oos.c_e12_m_min} to \metric{rev.w.s.oos.c_e12_m_max}~m, against \metric{rev.w.s.oos.c_all_m_min} to \metric{rev.w.s.oos.c_all_m_max}~m on all epochs; on the E3 units of calibration and test blocks it gives median area biases of \metric{rev.s.item4.oos.calibtest.c_e12.median_max_pct} to \metric{rev.s.item4.oos.calibtest.c_e12.median_min_pct}\% (\metric{rev.w.s.oos.validated_pct_max} to \metric{rev.w.s.oos.validated_pct_min}\% in the validated sectors), against \metric{rev.w.s.oos.c_all_pct_min} to \metric{rev.w.s.oos.c_all_pct_max}\% with the all-epoch shift on the same cells. Restricted to the band that emerged after 1995, the median intersection over union is \metric{rev.s.item4.overlap.shape.band_domain.iou.median_min}--\metric{rev.s.item4.overlap.shape.band_domain.iou.median_max} (\metric{rev.s.item4.overlap.shape.band_validated.iou.median_min}--\metric{rev.s.item4.overlap.shape.band_validated.iou.median_max} in its validated part), against \metric{rev.s.item4.overlap.shape.domain.iou.median_min}--\metric{rev.s.item4.overlap.shape.domain.iou.median_max} over the whole domain.

\emph{The deltaic test at sector level.} With three deltaic or vegetated sectors against four others there are \metric{rev.s.item5.sector.n_labellings} labellings, and with medians of three against four the most extreme labelling always ties with one other, so the smallest attainable one-sided $p$ is \metric{rev.s.item5.sector.min_attainable_p} (an erratum to Amendment 18 corrects the registered text, which gave a smaller value). For the three grids whose block test was significant the exact $p$ is \metric{rev.w.s.item5.sector_p_min}--\metric{rev.w.s.item5.sector_p_max}, and the rank-sum version gives \metric{rev.w.s.item5.sector_rank_p_min}--\metric{rev.w.s.item5.sector_rank_p_max}. Holm's adjustment over the ten registered block tests (five grids, two variants) leaves \metric{rev.s.item5.shape.block.n_significant_holm} shape-variant tests significant (adjusted $p$ \metric{rev.w.s.item5.block_p_holm_min}--\metric{rev.w.s.item5.block_p_holm_max}) and \metric{rev.s.item5.total.block.n_significant_holm} in the total variant. The comparison group holds \metric{rev.s.item5.shape.gebco2026.block.n_group} blocks against \metric{rev.s.item5.shape.gebco2026.block.n_rest} for GEBCO\_2026, with median distances per block of \metric{rev.s.item5.shape.gebco2026.block.group.distances_per_block_median} and \metric{rev.s.item5.shape.gebco2026.block.rest.distances_per_block_median}; the ninety-fifth percentile of its permutation null is \metric{rev.s.item5.shape.gebco2026.block.null_q95_km}~km.

\emph{Seabed change.} The registered check compares, at the levels crossed on both limbs, whether a grid fits the E1 or the later observed shoreline better: \metric{rev.s.item8.all.n} comparisons (five grids, seven sectors, four levels), of which E1 fits better in \metric{rev.s.item8.all.k_e1_better}. These are dependent; with one vote per sector and level the count is \metric{rev.s.item8.votes.k_e1_better} of \metric{rev.s.item8.votes.n} (two-sided binomial $p = \metric{rev.s.item8.votes.p}$). At the highest of the four levels the result reverses (E1 better in \metric{rev.s.item8.bin41.k_e1_better} of \metric{rev.s.item8.bin41.n}), where each comparison rests on few distances. The median distance between the E1 and later shorelines at the same level is \metric{rev.s.item8.hysteresis.west.median_km}~km in the west (within-epoch ninety-fifth percentile \metric{rev.s.item8.hysteresis.west.e2_halves_p95_km}~km), \metric{rev.s.item8.hysteresis.northeast.median_km}~km on the northeastern coast (\metric{rev.s.item8.hysteresis.northeast.e2_halves_p95_km}~km), \metric{rev.s.item8.hysteresis.ural.median_km}~km in the Ural delta (\metric{rev.s.item8.hysteresis.ural.e2_halves_p95_km}~km), \metric{rev.s.item8.hysteresis.volga.median_km}~km in the Volga delta, within its within-epoch scatter, and \metric{rev.s.item8.hysteresis.komsomolets.median_km}~km in Komsomolets Bay; in the interfluve and on the open shelf it is \metric{rev.s.item8.hysteresis.interfluve.median_km}~km. The check covers only the top of the record and says nothing about change between the source surveys and the Landsat era.

\emph{ETOPO 2022.} At all \metric{rev.g.etopo.lake_clamp.n_lake_cells} of its North Caspian lake-depth cells, ETOPO 2022 equals GEBCO\_2022 with values above \metric{rev.w.g.etopo.cap_m}~m set to \metric{rev.w.g.etopo.cap_m}~m (\metric{rev.g.etopo.lake_clamp.share_gebco_above_m29_pct}\% of the cells are so capped). Its documentation states that the lake layer takes GEBCO depths for very large lakes\citep{etopo2022_userguide,macferrin2025} but not the cap; the cap equals the Caspian surface elevation of HydroLAKES, whose outlines come foremost from the water-body data of February 2000\citep{hydrolakes_doc,messager2016}, and this link is our inference. In the band from \metric{h1.disp_band.edge_2_m} to \metric{h1.disp_band.edge_3_m}~m, \metric{rev.g.etopo.edge.m27p7_m27p0.lake_depth_edge.share_pct}\% of its predicted shoreline cells lie on the edge of that outline and \metric{rev.g.etopo.edge.m27p7_m27p0.fabdem.share_pct}\% on FABDEM. Against GEDI, its NMAD is \metric{rev.g.etopo.gedi.etopo2022.nmad_fabdem_m}~m at FABDEM points and \metric{rev.g.etopo.gedi.etopo2022.nmad_lake_depth_m}~m at lake-depth points; the block-bootstrap interval of the difference is \metric{rev.g.etopo.gedi.etopo2022.diff_lake_depth_minus_fabdem_ci_low_m} to \metric{rev.g.etopo.gedi.etopo2022.diff_lake_depth_minus_fabdem_ci_high_m}~m, and within the \metric{rev.g.etopo.gedi.within_block.n_blocks} blocks holding both sources the median difference is \metric{rev.g.etopo.gedi.within_block.etopo_median_diff_lake_minus_fabdem_m}~m, with FABDEM lower in \metric{rev.g.etopo.gedi.within_block.n_blocks_etopo_fabdem_lower} of them.

\emph{Great Salt Lake.} In the GEBCO type-identifier grids every band cell of both arms is coded as land, and SRTM15+ carries no sounding source there\citep{gebco2026,tozer2019}; ETOPO 2022 takes GEBCO depths only for the large lakes of its own list, which does not include Great Salt Lake, and elsewhere holds FABDEM's flattened water surfaces\citep{macferrin2025,copdem_handbook_i5}. The four SRTM-derived grids hold one flat height, \metric{rev.w.g.gsl.flat_height_navd88_m}~m (NAVD~88), on \metric{rev.w.g.gsl.flat_share_pct_min}--\metric{rev.w.g.gsl.flat_share_pct_max}\% of the band, \metric{rev.w.g.gsl.flat_above_feb2000_min_m}--\metric{rev.w.g.gsl.flat_above_feb2000_max_m}~m above the gauge levels of February 2000. At the control's levels the four SRTM-derived grids exceed the north-arm floor at \metric{rev.w.g.gsl.ctrl.north.four_exceeding_min} of \metric{rev.g.gsl.control_levels.gebco2026.shape.north.n_evaluable} levels and ETOPO 2022 at \metric{rev.g.gsl.control_levels.etopo2022.shape.north.n_exceeding}; in the south arm all five exceed it at \metric{rev.w.g.gsl.ctrl.south.exceeding_min}--\metric{rev.w.g.gsl.ctrl.south.exceeding_max} of \metric{rev.g.gsl.control_levels.gebco2026.shape.south.n_evaluable}. The registered control could not be retrieved: every scripted request met an HTTP \metric{rev.g.gsl.root2023.sciencebase_status} browser challenge, and the file exceeds the revision's download limit, so it was not run. The 3DEP tiles' versions served by Earth Engine are undocumented, so the legacy share of the south arm is a range.

\emph{H2.} The verdict is not supported in \metric{rev.g.h2.split.n_not_supported} of \metric{rev.g.h2.split.n_seeds} random reassignments of the test blocks (ATL08; the \metric{rev.g.h2.split.n_supported} others were run on ATL08 only, so they are at most partial support). No single test block reverses the NMAD difference (range \metric{rev.g.h2.influence.atl08.block.nmad.min_m} to \metric{rev.g.h2.influence.atl08.block.nmad.max_m}~m), but leaving out the northeastern sector, whose two test blocks hold \metric{rev.w.g.h2.ne_points_share_pct}\% of the points, reverses it for \metric{rev.g.h2.influence.atl08.sector.nmad.n_products_sign_change} grids, and the RMSE difference changes sign for all \metric{rev.g.h2.influence.atl08.block.rmse.n_products_sign_change} grids when one block is left out. Intervals for the interfluve and west, which have one test block each, are not estimable. On a common point set the grids' RMSE is \metric{rev.g.h2.blocksizes.atl08.50km.products.rmse_min_m}--\metric{rev.g.h2.blocksizes.atl08.50km.products.rmse_max_m}~m at \metric{h2.block_size_km}~km and \metric{rev.g.h2.blocksizes.atl08.12p5km.products.rmse_min_m}--\metric{rev.g.h2.blocksizes.atl08.12p5km.products.rmse_max_m}~m at \metric{params.block_size_sens_lo_km}~km, so the earlier rise at small blocks was a change of population. A dryness test from the observed extents instead of the elevation model agrees with it for \metric{rev.g.h2.dryness.atl08.agree.share_same_pct}\% of points and leaves H2 \metric{rev.g.h2.dryness.decision.matched}. Coverage of the nominal \metric{erl.text.conformal_pct}\% intervals is \metric{rev.g.h2.coverage.atl06.coverage_pct}\% (ATL06) and \metric{rev.g.h2.coverage.atl08.coverage_pct}\% (ATL08), with test-block bootstrap intervals of \metric{rev.g.h2.coverage.atl06.boot_low_pct}--\metric{rev.g.h2.coverage.atl06.boot_high_pct}\% and \metric{rev.g.h2.coverage.atl08.boot_low_pct}--\metric{rev.g.h2.coverage.atl08.boot_high_pct}\%.

\emph{Near-term emergence.} \emph{The grids as distributed} place their first metre partly on seabed that is already dry: \metric{rev.w.n.dry.unc1_pct_min}--\metric{rev.w.n.dry.unc1_pct_max}\% of each \metric{h3.scenario.1_m}~m band, and \metric{rev.w.n.dry.unc2_pct_min}--\metric{rev.w.n.dry.unc2_pct_max}\% of each \metric{h3.scenario.2_m}~m band, is not water in the 2025 and 2026 near-minimum extent, against \metric{rev.w.n.dry.shifted1_pct_min}--\metric{rev.w.n.dry.shifted1_pct_max}\% for the shifted grids. Their factor of \metric{rev.n.summary.uncorrected.1.ratio} at \metric{h3.scenario.1_m}~m comes from whole-metre terraces either side of the 2026 level, and a change of \metric{rev.w.n.sens.jitter_m}~m in that level reverses the ranking of GEBCO\_2022 and ETOPO 2022 (\metric{rev.n.sens.L26+0.3.uncorrected.gebco2022.1.domain_km2} against \metric{rev.n.sens.L26+0.3.uncorrected.etopo2022.1.domain_km2}~\kmsq{}). Without the EGM2008 realisation constant, their emergence at \metric{h3.scenario.0p5_m}~m rises from \metric{rev.n.summary.uncorrected.0p5.min_km2}--\metric{rev.n.summary.uncorrected.0p5.max_km2} to \metric{rev.n.sens.egm2008_constant_0.uncorrected.0p5.min_km2}--\metric{rev.n.sens.egm2008_constant_0.uncorrected.0p5.max_km2}~\kmsq{}.

\emph{Baselines.} The registered common baseline (a registered sensitivity) is the area of each grid's registered fused surface at $L_{2026}$, \metric{rev.n.baseline.sea_L2026.domain_min_km2}--\metric{rev.n.baseline.sea_L2026.domain_max_km2}~\kmsq{}. Cells not observed at the lowest audited level take the waterline model or the corrected grid there, so it marks as land seabed that the 2025 and 2026 imagery shows as water, and on its area version four grids hold more sea at $L_{2026}$ minus \metric{h3.scenario.1_m}~m than is observed at $L_{2026}$: $E(\Delta)$ at \metric{h3.scenario.1_m}~m is \metric{rev.n.summary.common_area.1.min_km2} to \metric{rev.n.summary.common_area.1.max_km2}~\kmsq{} (location version \metric{rev.n.summary.common_location.1.min_km2}--\metric{rev.n.summary.common_location.1.max_km2}~\kmsq{}). The Letter therefore uses a composite observed baseline, logged before it was computed: the observed sea at the lowest audited level, plus the 2025 and 2026 near-minimum water where that level was unobserved (\metric{rev.n.baseline.composite.baseline.domain_km2}~\kmsq{}). On it, $E(\Delta)$ at \metric{h3.scenario.1_m}~m is \metric{rev.n.composite.summary.common_location.1.min_km2}--\metric{rev.n.composite.summary.common_location.1.max_km2}~\kmsq{}, and \metric{rev.n.composite.sens.L26+0.3.common_location.1.min_km2}--\metric{rev.n.composite.sens.L26+0.3.common_location.1.max_km2}~\kmsq{} or \metric{rev.n.composite.sens.L26-0.3.common_location.1.min_km2}--\metric{rev.n.composite.sens.L26-0.3.common_location.1.max_km2}~\kmsq{} with $L_{2026}$ raised or lowered by \metric{rev.w.n.sens.jitter_m}~m. No meaningful part of the first metre is observed directly: without the record-low rule, \metric{rev.n.observed_band_v2.no_envelope.band_13_12.domain_km2}~\kmsq{} lies between the two lowest audited levels.

\emph{Envelopes.} The registered scale-aware rule gives no finite envelope, because the E1 calibration pairs show no observed change, so their scores are infinite. With the E3 pairs alone (the registered lowest-epoch sensitivity) the quantile is \metric{rev.w.n.env.q_e3_shifted_min}--\metric{rev.w.n.env.q_e3_shifted_max} on the shifted grids and the upper bounds remain infinite. Normalised by the decline (a logged deviation), the envelopes at \metric{h3.scenario.1_m}~m span \metric{rev.w.n.env.shifted1.lo_min}--\metric{rev.w.n.env.shifted1.hi_max}~\kmsq{} on the shifted grids and \metric{rev.w.n.env.composite1.lo_min}--\metric{rev.w.n.env.composite1.hi_max}~\kmsq{} from the observed shoreline; declines beyond \metric{rev.n.env.max_pair_decline_m}~m, including the \metric{h3.scenario.3_m}~m and \metric{h3.scenario.5_m}~m scenarios, are uncalibrated. The paired difference between the registered correction and the uncorrected grids lies outside its envelope in \metric{rev.n.paired.linear_hold.n_outside} of \metric{rev.n.paired.linear_hold.n_cases} grid-by-decline cases (\metric{rev.n.paired.constant.n_outside} for the constant form); the registered verdict stands.

\emph{Vertical-error simulation.} Fields at each grid's ICESat-2 dispersion were simulated with correlation ranges of \metric{rev.w.n.sim.range_primary_km}~km and \metric{rev.w.n.sim.range_secondary_km}~km. The shifted grids' $E(\Delta)$ at \metric{h3.scenario.1_m}~m has a standard deviation of \metric{rev.w.n.sim.r150km.sd_shifted1_km2_min}--\metric{rev.w.n.sim.r150km.sd_shifted1_km2_max}~\kmsq{} (\metric{rev.w.n.sim.r150km.sd_rel_shifted1_pct_min}--\metric{rev.w.n.sim.r150km.sd_rel_shifted1_pct_max}\%) at \metric{rev.w.n.sim.range_primary_km}~km and \metric{rev.w.n.sim.r50km.sd_shifted1_km2_min}--\metric{rev.w.n.sim.r50km.sd_shifted1_km2_max}~\kmsq{} (\metric{rev.w.n.sim.r50km.sd_rel_shifted1_pct_min}--\metric{rev.w.n.sim.r50km.sd_rel_shifted1_pct_max}\%) at \metric{rev.w.n.sim.range_secondary_km}~km. Five perturbed copies of one grid agree on \metric{rev.w.n.sim.r150km.agree_unc1_ref_p05_pct_min}--\metric{rev.w.n.sim.r150km.agree_unc1_ref_p95_pct_max}\% of their \metric{h3.scenario.1_m}~m band at \metric{rev.w.n.sim.range_primary_km}~km and \metric{rev.w.n.sim.r50km.agree_unc1_ref_p05_pct_min}--\metric{rev.w.n.sim.r50km.agree_unc1_ref_p95_pct_max}\% at \metric{rev.w.n.sim.range_secondary_km}~km. The five grids agree on \metric{rev.w.n.sim.agree_unc1_observed_pct}\%, at or above \metric{rev.w.n.sim.r150km.agree_unc1_groups_below_pct_min}--\metric{rev.w.n.sim.r150km.agree_unc1_groups_below_pct_max}\% and \metric{rev.w.n.sim.r50km.agree_unc1_groups_below_pct_min}--\metric{rev.w.n.sim.r50km.agree_unc1_groups_below_pct_max}\% of the reference groups at the two ranges, so their low agreement is not evidence of error beyond the validated vertical error. At \metric{h3.scenario.2_m}~m the grids agree more than independent copies do at both ranges, as their shared soundings lead one to expect.

\emph{Dry seabed and isolated water.} Water cut off from the sea is \metric{rev.w.n.isolated.unc1_pct_min}--\metric{rev.w.n.isolated.unc1_pct_max}\% of the uncorrected \metric{h3.scenario.1_m}~m bands, \metric{rev.w.n.isolated.shifted1_pct_min}--\metric{rev.w.n.isolated.shifted1_pct_max}\% of the shifted ones and at most \metric{rev.w.n.isolated.composite1_pct_max}\% from the observed shoreline; nearly all emergence is dry seabed.

\emph{Exposure.} Protected areas hold \metric{rev.w.n.wdpa.composite.with_pct_min}--\metric{rev.w.n.wdpa.composite.with_pct_max}\% of the first metre from the observed shoreline and \metric{rev.w.n.wdpa.shifted.with_pct_min}--\metric{rev.w.n.wdpa.shifted.with_pct_max}\% on the shifted grids, but \metric{rev.w.n.wdpa.composite.without_pct_min}--\metric{rev.w.n.wdpa.composite.without_pct_max}\% and \metric{rev.w.n.wdpa.shifted.without_pct_min}--\metric{rev.w.n.wdpa.shifted.without_pct_max}\% without the database record of the Kazakh State Reserved Zone, drawn an order of magnitude larger than its reported area. The Volga Delta Ramsar site holds \metric{rev.n.exp.wdpa_ramsar.common_location_composite.1.min_km2}--\metric{rev.n.exp.wdpa_ramsar.common_location_composite.1.max_km2}~\kmsq{} of the first metre from the observed shoreline and \metric{rev.n.exp.wdpa_ramsar.shifted.1.min_km2}--\metric{rev.n.exp.wdpa_ramsar.shifted.1.max_km2}~\kmsq{} on the shifted grids. Almost all OpenStreetMap features tagged as ports are piers or ferry landings, many far up the deltas, so port counts are not reported; the Ural Delta Ramsar site is a point record and is not estimated. The location envelope of the exposure counts was not run.

\emph{Pathways.} Across the primary pathways and the five grids, the pathway carries \metric{rev.w.n.var.pathway_pct_min}--\metric{rev.w.n.var.pathway_pct_max}\% of the sum of squares of emerged area at 2035 and 2050; from the observed shoreline in 2035 the grids carry \metric{rev.w.n.var.composite2035.grid_area_pct}\% (\metric{rev.w.n.var.composite2035.grid_log_pct}\% on a log scale). Calendar time is not re-anchored: each pathway's change after its base year is added to the observed 2026 level, which drops any mismatch with the observed fall of \metric{rev.w.n.anchor.observed_fall_m}~m between 2021 and 2026, against \metric{rev.n.anchor.primary.min_fall_2021_2026_m} to \metric{rev.n.anchor.primary.max_fall_2021_2026_m}~m in the primary pathways. The observed-trend pathway depends on its window: at \metric{rev.w.n.trend.1996_2021.fall_cm_per_yr}~cm per year (1996--2021) it reaches \metric{h3.scenario.1_m}~m in \metric{rev.w.n.trend.1996_2021.year_1m}, at \metric{rev.w.n.trend.2006_2025.fall_cm_per_yr}~cm per year (2006--2025) in \metric{rev.w.n.trend.2006_2025.year_1m}, and at \metric{rev.w.n.trend.2016_2025.fall_cm_per_yr}~cm per year (2016--2025) in \metric{rev.w.n.trend.2016_2025.year_1m}. The range pathways are in table~\ref{tab:ed_scenario_dates}.
